\documentclass[11pt]{article}
\usepackage{times}
\usepackage{fullpage}
\usepackage{url}
\usepackage{hyperref}
\usepackage{fancyhdr}
\usepackage{graphicx}
\usepackage{enumitem}
\usepackage{subcaption}
\usepackage{amsmath, amsfonts, amsthm, fouriernc}
\usepackage{IEEEtrantools}
\usepackage{parskip}
\usepackage{booktabs}

\pagestyle{fancy}
\addtolength{\headheight}{2em}
\addtolength{\headsep}{1.5em}
\lhead{\large{ME 2801}}
\rhead{\large{Solutions}}

\usepackage{sectsty}
\sectionfont{\fontsize{12}{8}\selectfont}

% Auto-numbered exercise headings: \exercise{Title} -> "Exercise N --- Title"
\newcounter{exercise}
\newcommand{\exercise}[1]{%
  \refstepcounter{exercise}%
  \section*{Exercise \theexercise\ --- #1}%
}

\begin{document}

\begin{center}
  \Large{\bf{Solutions: Steady-State Error}}
\end{center}

Numerical verification for every exercise is in
\texttt{assignment\_solutions.m}.  

\textbf{Notation.}  Throughout these solutions:
\begin{itemize}
  \item $G_{ol}(s)$ --- the open-loop (forward-path) transfer function.
  \item $G_{cl}(s)$ --- the closed-loop transfer function.
  \item $G(s)$ --- the transfer function appearing inside the SSE master
    formula $e(\infty) = \lim_{s \to 0} sR(s)/(1 + G(s))$.  In our
    unity-feedback setup, $G(s) = G_{ol}(s)$.
\end{itemize}

% ---------------------------------------------------------------
\exercise{Three Test Inputs}

\textbf{(a)}  $G_{ol}(s)$ has no factors of $s$ in the denominator
$\Rightarrow$ \textbf{Type 0}.

\textbf{(b)}  Steady-state errors follow from the master formula
\[
  e(\infty) = \lim_{s \to 0} \frac{s\,R(s)}{1 + G(s)},
\]
with each input $R(s) = 1/s,\ 1/s^2,\ 1/s^3$ in turn.  The relevant DC
limits of $G(s) = G_{ol}(s)$ are
\[
  K_x \equiv \lim_{s \to 0} G(s) = \frac{600}{(10)(25)} = 2.4,
  \qquad
  \lim_{s \to 0} s\,G(s) = 0,
  \qquad
  \lim_{s \to 0} s^2\,G(s) = 0,
\]
giving
\[
  e_\mathrm{step} = \frac{1}{1 + K_x} = \frac{1}{3.4} \approx \boxed{0.294},
  \qquad
  e_\mathrm{ramp} = \boxed{\infty},
  \qquad
  e_\mathrm{para} = \boxed{\infty}.
\]

\textbf{(c)}  In MATLAB:
\begin{itemize}
  \item \emph{Calculate.}  \texttt{G\_cl = feedback(G\_ol,1); 1 - dcgain(G\_cl)}
    returns 0.2941.
  \item \emph{Simulate.}  \texttt{step(G\_cl)} settles at
    $0.706 = 1 - 0.294$, so the error signal $r - y$ settles at $0.294$.
\end{itemize}
Both match the analytical answer.

\textbf{(d)}  $G_{ol}(s)$ contains \emph{no integrators} (no factor of $s$
in the denominator).  Without integration in the forward path, the
controller cannot accumulate the unbounded bias required to keep up with
an unbounded reference; the error grows without limit for the ramp and
parabola.

% ---------------------------------------------------------------
\exercise{Design a P-Only Velocity Tracking Compensator}

\textbf{(a)}  The combined open-loop (forward path) is
\[
  G_{ol}(s) = C(s)\,G_p(s) = \frac{K_p\,(s+5)}{s\,(s+8)(s+12)},
\]
which has one factor of $s$ in the denominator $\Rightarrow$ \textbf{Type 1}.

\textbf{(b)}  For a unit ramp input, $R(s) = 1/s^2$, so the master formula
gives
\[
  e_\mathrm{ramp} = \lim_{s \to 0} \frac{s\,R(s)}{1 + G(s)}
                  = \lim_{s \to 0} \frac{1/s}{1 + G(s)}.
\]
With $G(s) = G_{ol}(s)$, near the origin the integrator dominates,
\[
  G(s) = \frac{K_p\,(s+5)}{s\,(s+8)(s+12)}
       \;\xrightarrow{s \to 0}\;
       \frac{5\,K_p}{96\,s},
\]
which is large compared to $1$, so
\[
  e_\mathrm{ramp} = \lim_{s \to 0} \frac{1/s}{5\,K_p/(96\,s)}
                  = \frac{96}{5\,K_p}.
\]
A 4\,\% ramp error requires $96/(5\,K_p) = 0.04$:
\[
  \boxed{K_p = \frac{96}{5 \cdot 0.04} = 480}.
\]

\textbf{(c)}  In the script, \texttt{G\_cl = feedback(Kp*G\_p, 1)};
\texttt{pole(G\_cl)} with $K_p = 480$ returns roots near
$-4.46 \pm 14.2j$ and $-11.1$ --- all in the left half-plane, so the closed
loop is stable.  This is what makes the SSE result in (b) meaningful: the
chapter's formula assumes a stable closed loop.

% ---------------------------------------------------------------
\exercise{USV Surge: Predict Before You Simulate}

Let $C(s)$ denote the controller (P or PI), $G_p(s) = 5.1/(1.3s+1)$ the
plant, and $K_{\!f\!f} = 1/G_p(0) = 1/5.1$ the feedforward gain.  With
feedforward the closed-loop transfer function is
\[
  G_{cl}(s) = \frac{G_p(s)\bigl(C(s) + K_{\!f\!f}\bigr)}{1 + C(s)\,G_p(s)},
\]
and the step error is $e_\mathrm{ss} = 1 - G_{cl}(0)$.

\textbf{(a) P only.}  $C(s) = K_P = 0.5$.  The forward path
$G_{ol}(s) = K_P\,G_p(s)$ is Type~0 with $K_x = K_P \cdot G_p(0) = 0.5
\cdot 5.1 = 2.55$:
\[
  e_\mathrm{ss} = \frac{1}{1 + K_x} = \frac{1}{3.55}
                \approx \boxed{0.282\;\mathrm{m/s}}.
\]

\textbf{(b) PI.}  Adding $K_I/s$ to the controller introduces an integrator
in $G_{ol}(s)$ $\Rightarrow$ Type~1, so
\[
  e_\mathrm{ss} = \boxed{0}.
\]

\textbf{(c) P + FF.}  $C(0) = 0.5$, so
\[
  G_{cl}(0) = \frac{5.1\,(0.5 + 1/5.1)}{1 + 5.1 \cdot 0.5}
            = \frac{5.1 \cdot 0.696}{3.55}
            = \frac{3.55}{3.55}
            = 1
  \;\Longrightarrow\;
  e_\mathrm{ss} = \boxed{0}.
\]

\textbf{(d) PI + FF.}  $C(s) \to \infty$ as $s \to 0$ (the $K_I/s$ term
dominates), so $G_{cl}(0) = 1$ as in (c) and
\[
  e_\mathrm{ss} = \boxed{0}.
\]

\textbf{(e)}  At steady state in a stable closed loop, every signal is
constant, so $\dot{u}_I(t) = 0$.  But $\dot{u}_I = K_I\,e(t)$, so
$\dot{u}_I = 0$ forces $e = 0$ (assuming $K_I \neq 0$).  The integrator
\emph{adapts} until the error has been driven to zero --- it does not need
to know $K_\mathrm{plant}$ to do so.

% ---------------------------------------------------------------
\exercise{USV Yaw: Ramp Tracking Design}

\textbf{(a)}  Heading is the integral of yaw rate, so the heading plant
picks up an extra factor of $1/s$ relative to the yaw-rate plant:
\[
  G_r(s) = \frac{1}{s}\,G_\mathrm{yaw\text{-}rate}(s)
         = \frac{0.82}{s\,(0.61\,s + 1)}.
\]

\textbf{(b)}  The combined open-loop is
\[
  G_{ol}(s) = K_r\,G_r(s) = \frac{0.82\,K_r}{s\,(0.61\,s + 1)},
\]
which has one factor of $s$ in the denominator $\Rightarrow$ \textbf{Type 1}.
($K_r$ is just a scalar gain; it cannot add or remove integrators.)

\textbf{(c)}  The reference is a ramp of slope $5^\circ$/s, so $R(s) =
5/s^2$.  The master formula gives
\[
  e_\mathrm{ss} = \lim_{s \to 0} \frac{s\,R(s)}{1 + G(s)}
                = \lim_{s \to 0} \frac{5/s}{1 + G(s)}.
\]
With $G(s) = G_{ol}(s)$, near the origin the integrator dominates,
\[
  G(s) = \frac{0.82\,K_r}{s\,(0.61\,s + 1)}
       \;\xrightarrow{s \to 0}\;
       \frac{0.82\,K_r}{s},
\]
which is large compared to $1$, so
\[
  e_\mathrm{ss} = \lim_{s \to 0} \frac{5/s}{0.82\,K_r/s}
                = \frac{5}{0.82\,K_r} \;\;[\text{degrees}].
\]

\textbf{(d)}  Spec: $|e_\mathrm{ss}| \le 1^\circ$:
\[
  \frac{5}{0.82\,K_r} \le 1
  \;\Longrightarrow\;
  K_r \ge \frac{5}{0.82}
        \approx \boxed{6.10}.
\]

\textbf{(e)}  In MATLAB, \texttt{lsim} on the closed loop with $K_r = 6.10$
and a $5t$ reference shows the output settling parallel to the reference
with a constant error of $1.0^\circ$, confirming the spec.

\end{document}
